The goal is to adapt a pretrained VLA to complete language-conditioned manipulation tasks on a target robot without collecting human demonstrations. We study multi-stage pick-and-place tasks: an instruction such as \emph{put the apple on the plate} requires approach, grasp, lift, transport, and release. The policy receives camera images, proprioception, and an instruction, predicts an action chunk, and executes a prefix before observing again.

During acquisition, each skill has its own instruction, success condition, and starting states. During composition and final evaluation, the robot starts with an empty gripper and receives one complete-task instruction. This distinction prevents isolated skill performance from being treated as full-task success. Algorithm~\ref{alg:paper_pipeline} summarizes the pipeline: acquire simple skills, consolidate them, and learn the complete instruction.

We use SmolVLA as the pretrained model~\cite{smolvla}. On the CDPR and WidowX, we map its action prior to three translation commands, yaw, and gripper opening. A newly initialized residual MLP receives the mapped prior and observation features and predicts corrections to an eight-action chunk. The robot executes four actions before observing again. Earlier training also adapts the pretrained expert with low-rank adapters. The final stage freezes the expert and its adapters and continues training the retained residual controller.
